\section{Reconstruction of the attainable volume profile}
\label{sec:operational-reconstruction}
The preceding theorems determine a prediction law from an attainable
flag. There is a converse: the entire integer-budget regret curve
recovers every product in that flag, up to uniform constants. Thus the
truncated volume profile is also necessary for uniform comparison of
prediction laws. We first prove the elementary envelope identity on
which this converse rests. Its role is to extract the invariant from
the two experiment-specific classifications; it is not a replacement
for their attainment, global covering, or causal arguments.

\subsection{An envelope identity with integer budgets}
For an ordered list $s_1\ge\cdots\ge s_p\ge0$, put
\[
 V_0(s)=1,\qquad V_\ell(s)=\prod_{i=1}^{\ell}s_i,
 \qquad
 e_s(b)=\max_{1\le j\le p}\left(\frac{V_j(s)}b\right)^{1/j}
 \quad(b\ge1).
\]
Zero entries and their products are retained. The extension to real
$b$ is an analytic device; an actual state alphabet always has an
integer cardinality.

\begin{lemma}[Recovery of all initial products]
\label{lem:integer-envelope-duality}
For every $1\le\ell\le p$,
\begin{equation}\label{eq:continuous-envelope-recovery}
 V_\ell(s)=\inf_{b\ge1}b\,e_s(b)^\ell.
\end{equation}
For integer budgets one has
\begin{equation}\label{eq:integer-envelope-recovery}
 V_\ell(s)\le
 \inf_{M\in\N,\,M\ge1}M e_s(M)^\ell
 \le2V_\ell(s).
\end{equation}
Both assertions include zero products. If $s_\ell>0$, equality in
\eqref{eq:continuous-envelope-recovery} holds at
$b_\ell=V_\ell(s)/s_\ell^\ell\ge1$.
\end{lemma}
\begin{proof}
The $\ell$th term in the maximum gives
$b e_s(b)^\ell\ge V_\ell(s)$. Suppose first that $s_\ell>0$.
For $j<\ell$, monotonicity of the list gives
\[
 \frac{V_\ell}{V_j}=\prod_{i=j+1}^{\ell}s_i
       \ge s_\ell^{\ell-j},
\]
and hence $V_j/b_\ell\le s_\ell^j$. For $j>\ell$, the reverse
product estimate $V_j/V_\ell\le s_\ell^{j-\ell}$ gives the same
inequality. The $\ell$th term is equal to $s_\ell$, so
$e_s(b_\ell)=s_\ell$. Taking $M_\ell=\lceil b_\ell\rceil$ and
using that $e_s$ is nonincreasing, we obtain
\[
 M_\ell e_s(M_\ell)^\ell
 \le\lceil b_\ell\rceil s_\ell^\ell
 \le2b_\ell s_\ell^\ell=2V_\ell.
\]
This proves both assertions for positive products without rounding a
real budget down to an inadmissible alphabet size.

Let $r$ be the number of positive entries. If $r=0$, every quantity
in the assertions is zero. If $0<r<\ell$, then for
$b\ge b_r=V_r/s_r^r$ the last positive term is maximal. Indeed,
with $t=(V_r/b)^{1/r}\le s_r$, the inequality
$V_r/V_j\ge s_r^{r-j}\ge t^{r-j}$ gives $V_j/b\le t^j$ for
$j<r$. Consequently
\[
 b e_s(b)^\ell=V_r^{\ell/r}b^{1-\ell/r}\longrightarrow0.
\]
The same limit holds along the integers. Since $V_\ell=0$, this
settles every remaining case.
\end{proof}

The equality uses the ordered-product structure. It would be incorrect
to apply it to an arbitrary list of coefficients in a maximum of power
laws: some such coefficients can lie strictly below the envelope. Here
$\log V_\ell$ is concave on the positive indices, and each product
has the supporting budget exhibited in the proof. Repeated scales may
share that budget; their products are still recovered.

\subsection{A complete invariant for uniform regret comparison}
For any nonnegative curve $R:\N\to[0,\infty)$ define
\begin{equation}\label{eq:operational-volume-transform}
 \mathcal I_0(R)=1,\qquad
 \mathcal I_\ell(R)=\inf_{M\in\N,\,M\ge1}
                       M R(M)^{\ell/2}\quad(\ell\ge1).
\end{equation}
The infimum refers to the whole optimal-regret curve, not to finitely
many observations of an implemented filter. No numerical inversion
procedure or finite-budget estimator is implicit in this definition.

\begin{theorem}[Operational reconstruction and comparison]
\label{thm:operational-reconstruction}
Suppose a family of checkpoint regret curves satisfies
\begin{equation}\label{eq:reconstruction-hypothesis}
 c e_s(M)^2\le R_s(M)\le C e_s(M)^2
 \qquad(M\in\N,\ M\ge1)
\end{equation}
with $c,C>0$ uniform in the family and ordered lists of a fixed
maximum length $p$, padded with zeros when necessary. Then
\begin{equation}\label{eq:risk-recovers-products}
 c^{\ell/2}V_\ell(s)\le\mathcal I_\ell(R_s)
             \le2C^{\ell/2}V_\ell(s),\qquad 1\le\ell\le p.
\end{equation}
In particular the positive products, the scale orders of their successive
ratios, and their uniform orders of degeneration are
invariants of the regret curve up to comparison constants.

For two such families, their curves are uniformly comparable for all
integer $M\ge1$ if and only if all their corresponding products
$V_\ell$ are uniformly comparable, including their zero sets.
The comparison concerns finite-resolution prediction laws, not an
isomorphism of experiments or a reconstruction of their hidden parameters.
\end{theorem}
\begin{proof}
Raise \eqref{eq:reconstruction-hypothesis} to the power $\ell/2$,
multiply by $M$, and take the infimum. Lemma~\ref{lem:integer-envelope-duality}
gives \eqref{eq:risk-recovers-products}. For a positive product the
ratio $\mathcal I_\ell/\mathcal I_{\ell-1}$ is comparable to
$V_\ell/V_{\ell-1}=s_\ell$, with $\mathcal I_0=V_0=1$.
The first zero product and all its successors are detected as exact
zeros of the transform; no quotient $0/0$ is formed.

Suppose $A^{-1}R_s(M)\le\widetilde R_t(M)\le A R_s(M)$ uniformly.
The transform preserves these inequalities with constants $A^{\ell/2}$.
Applying \eqref{eq:risk-recovers-products} to both curves proves uniform
comparison of their products. Conversely, if
$B^{-1}V_\ell(s)\le V_\ell(t)\le B V_\ell(s)$ for all $\ell$,
with $B\ge1$, then
\[
 B^{-2}e_s(M)^2\le e_t(M)^2\le B^2e_s(M)^2.
\]
Indeed each branch has comparison factor $B^{2/\ell}\le B^2$;
taking maxima preserves the bounds. The two instances of
\eqref{eq:reconstruction-hypothesis} give the asserted comparison of
curves. Every constant depends only on the stated uniform constants
and the fixed maximum length, not on a positive lower bound for any
individual scale.
\end{proof}

For positive rank $r$, the proof of the lemma also gives
\[
 \lim_{M\to\infty}\frac{\log R_s(M)}{\log M}=-\frac2r.
\]
Thus the usual dimension slope is one consequence of the curve; the
products recovered in \eqref{eq:risk-recovers-products} give its
additional finite-resolution information. Uniformly equivalent norms
on the same prediction vector change every loss, and hence every
optimal curve, by uniform constants. The recovered product orders are
therefore independent of a uniformly conditioned choice of prediction
coordinates. Parameter-dependent changes with unbounded condition
number are not covered by this observation.

\subsection{The collision and contrast invariants}
For a monomial checkpoint define the past-truncated products
\[
 \widehat{\mathcal V}_{n,m;\ell}(a)=
 \begin{cases}
  \mathcal V_{m,\ell}(a),&1\le\ell\le p_{n,m},\\
  0,&\ell>p_{n,m}.
 \end{cases}
\]
The truncation is part of the invariant, not a convention for suppressing
unattained directions.

\begin{corollary}[Recovery at every additive collision]
\label{cor:monomial-operational-recovery}
Under the hypotheses of Theorem~\ref{thm:intrinsic-checkpoint}, let
$R^a(M)$ be either $\overline R^a_{M;n,m}$ or
$R^{a,\max}_{M;n,m}$. For each fixed finite maximum index,
\begin{equation}\label{eq:monomial-operational-recovery}
 \mathcal I_\ell(R^a)\asymp
       \widehat{\mathcal V}_{n,m;\ell}(a),\qquad \ell\ge1,
\end{equation}
uniformly on $K$, including all additive collision strata. Consequently
these truncated volumes are a complete invariant for uniform comparison
of the corresponding checkpoint regret curves.
\end{corollary}
\begin{proof}
Take $s_j=d_j$ for $1\le j\le p_{n,m}$ from
Lemma~\ref{lem:leja-scales} and pad by zeros. The same lemma gives
$V_\ell(s)\le\mathcal V_{m,\ell}\le\ell!V_\ell(s)$ through the
truncation. These factors are uniform at fixed horizon, including when
both sides vanish. Theorem~\ref{thm:intrinsic-checkpoint} therefore
verifies \eqref{eq:reconstruction-hypothesis} for either regret curve.
Theorem~\ref{thm:operational-reconstruction} proves the assertion,
including zero products beyond the truncation.
\end{proof}

The conclusion recovers the observable volume orders, not the complete
arrangement of the future exponents. In particular it does not recover
directions inaccessible to this past. The latter may still occur in the
prior-ambiguity body of Section~\ref{sec:directional-geometry}, whose
information convention is different.

\begin{corollary}[Recovery of paired acquisition--observation scales]
\label{cor:circular-operational-recovery}
For either checkpoint regret curve in
Theorem~\ref{thm:circle-resolution}, with $k=\min(n,m)$, uniformly on
its fixed-horizon contrast interval,
\begin{equation}\label{eq:circle-operational-recovery}
 \begin{split}
 \mathcal I_{2j-1}(R^\tau)&\asymp\tau^{2j^2},\\
 \mathcal I_{2j}(R^\tau)&\asymp\tau^{2j(j+1)},
             \qquad 1\le j\le k.
 \end{split}
\end{equation}
The transforms vanish at every higher index. For $\tau>0$ the
successive nonzero ratios recover both copies of each scale
$\tau^{2j}$; at $\tau=0$ all positive-order transforms vanish.
\end{corollary}
\begin{proof}
Use the real list
$(\tau^2,\tau^2,\tau^4,\tau^4,\ldots,\tau^{2k},\tau^{2k})$.
Its products of lengths $2j-1$ and $2j$ have respectively the exponents
$2j^2$ and $2j(j+1)$. The proof of
Theorem~\ref{thm:circle-resolution} identifies its full envelope with
$\mathcal H_{n,m}$, including the odd products even though they add no
larger branch. Theorem~\ref{thm:operational-reconstruction} now applies.
Ratios of consecutive products have exponent $2j$ on both members of
the $j$th pair. Zero contrast and indices above $2k$ are the zero-product
cases of Lemma~\ref{lem:integer-envelope-duality}.
\end{proof}

These reconstructions concern checkpoint curves with the experiment
known. Taking the maximum over checkpoints produces the causal laws
already proved, but need not retain the individual checkpoint curves.
Nor is a curve with a positive common-name input-error floor an instance
of \eqref{eq:reconstruction-hypothesis}. The distinct uncertainty law in
Theorem~\ref{thm:sharp-common-moments} and all of its hypotheses remain
in force.
